\begingroup
\renewcommand{\arraystretch}{1.45}
\fancyhead[L]{\small ARITHMETIC–GEOMETRIC MEAN.}
\Needspace{36\baselineskip}
\AGMArticle{15}
[The second-order differential equation for the terms of the arithmetic–geometric mean progression takes a form containing differentials only of quotients of the variables if, with an arbitrary quantity $e$, one forms the expression
$D+\mathrm d(\frac{1}{\Delta}\,\allowbreak \mathrm d\log e)+\frac{1}{\Delta}(\mathrm d\log e)^2$;
namely, this becomes
\Needspace{10\baselineskip}
\[
\begin{aligned}&-\surd(eea^{\mathrm n}b^{\mathrm n})\mathbin{.}\mathrm d
\left(\frac{1}{\Delta a^{\mathrm n}b^{\mathrm n}}\mathbin{.}\mathrm d
\surd\frac{a^{\mathrm n}b^{\mathrm n}}{ee}\right)\\&\quad-\frac{1}{4\Delta}\left(\mathrm d\log\frac{a^{\mathrm n}}{b^{\mathrm n}}\right)^2\end{aligned}
\]
an expression whose value is therefore independent of $n$ and also does not change if $b^{\mathrm n}$ and $c^{\mathrm n}$, or $c^{\mathrm n}$ and
$a^{\mathrm n}$, are substituted for $a^{\mathrm n}$ and $b^{\mathrm n}$. For steadily increasing positive values and for negative values of $n$, this expression is transformed into
\Needspace{10\baselineskip}
\[
\begin{gathered}-e\mathrm M(a,b)\mathbin{.}\mathrm d
\left\{\frac{1}{\Delta\mathrm M(a,b)^2}\,\mathrm d\frac{\mathrm M(a,b)}{e}\right\}
\\[7pt]\mbox{and into}\\[7pt]
-e\mathrm M(a,c)\mathbin{.}\mathrm d
\left\{\frac{1}{\Delta\mathrm M(a,c)^2}\,\mathrm d\frac{\mathrm M(a,c)}{e}\right\}\end{gathered}
\]
whereas, for $n$ equal to zero and for $a$, $b$, $c$ as special values of $e$, it becomes respectively $+\Delta bbcc,\quad-\Delta ccaa,\quad-\Delta aabb.$

\Needspace{24\baselineskip}
\par Thus, writing for brevity
\Needspace{10\baselineskip}
\[
\begin{gathered}\surd\frac{a}{\mathrm M(a,b)}=p,\\[7pt]
\surd\frac{b}{\mathrm M(a,b)}=q,\\[7pt]
\surd\frac{c}{\mathrm M(a,b)}=r,\\[7pt]
-\pi\frac{\mathrm M(a,b)}{\mathrm M(a,c)}=\log y\end{gathered}
\]
\Needspace{25\baselineskip}
and taking account of the equation for $\Delta$ in Article 13 gives:
\Needspace{23\baselineskip}
\[
\begin{aligned}
\displaystyle\frac{\Delta}{2}\mathrm M(a,b)^2
&=\displaystyle\frac14\mathrm d\log y\\&=\frac{1}{p^4}\,\mathrm d\log\frac{r}{q}\\&=\frac{1}{q^4}\,\mathrm d\log\frac{r}{p}\\&=\frac{1}{r^4}\,\mathrm d\log\frac{p}{q}\\[8pt]
&=\displaystyle-\frac{pp}{q^4r^4}\,\mathrm d
\left(\frac{1}{\mathrm d\log y}\mathbin{.}\mathrm d\frac{1}{pp}\right)\\&=\frac{qq}{r^4p^4}\,\mathrm d
\left(\frac{1}{\mathrm d\log y}\mathbin{.}\mathrm d\frac{1}{qq}\right)\\&=\frac{rr}{p^4q^4}\,\mathrm d
\left(\frac{1}{\mathrm d\log y}\mathbin{.}\mathrm d\frac{1}{rr}\right)
\end{aligned}
\]
\Needspace{18\baselineskip}
and the expressions for $-\frac{\Delta}{2}\mathrm M(a,c)^2$ in terms of
\Needspace{10\baselineskip}
\[
\begin{gathered}\surd\frac{a}{\mathrm M(a,c)},\\[7pt]
\surd\frac{c}{\mathrm M(a,c)},\\[7pt]
\surd\frac{b}{\mathrm M(a,c)},\\[7pt]
-\pi\frac{\mathrm M(a,c)}{\mathrm M(a,b)}\end{gathered}
\]
have the same form.
\Needspace{25\baselineskip}
\par Eliminating any two of the three quantities $p,q,r$ gives
\Needspace{10\baselineskip}
\[
\begin{aligned}&\left\{\frac{1}{p^4}\frac{1}{\mathrm d\log y}\,\mathrm d\log
\left[\frac{16}{p^6}\frac{1}{\mathrm d\log y}\,\mathrm d
\left(\frac{1}{\mathrm d\log y}\mathbin{.}\mathrm d\frac{1}{pp}\right)\right]\right\}^{2}\\&\quad-\frac{16}{p^6}\frac{1}{\mathrm d\log y}\,\mathrm d
\left(\frac{1}{\mathrm d\log y}\mathbin{.}\mathrm d\frac{1}{pp}\right)-1=0\end{aligned}
\]
\SourceLeaf{392}%
as a differential equation for $p$ as well as for $q$ and for $r$, that is, for
$\surd{\frac{a}{\mathrm M(a,b)}}$, $\surd{\frac{b}{\mathrm M(a,b)}}$ and
$\surd{\frac{c}{\mathrm M(a,b)}}$; and likewise, if
$-\pi\frac{\mathrm M(a,c)}{\mathrm M(a,b)}=\log z$ is substituted in it for $\log y$,
as a differential equation for $\surd{\frac{a}{\mathrm M(a,c)}}$,
$\surd{\frac{c}{\mathrm M(a,c)}}$ and $\surd{\frac{b}{\mathrm M(a,c)}}$.]

\Needspace{28\baselineskip}
\AGMArticle{16}
[The quotients of the quantities $a,b,c,\mathrm M(a,b),\mathrm M(a,c)$ can be represented by series proceeding in powers of $y^{\frac14}$ or $z^{\frac14}$ with the aid of the fundamental theorems of the algorithm to be investigated here, for example, in the following way.

By the definition of $y$ and the theorem stated in Article 12 for the algorithm extended backwards, the following holds:
\Needspace{10\baselineskip}
\[
-\pi\frac{\mathrm M(a,b)}{\mathrm M(a,c)}
=-\pi\frac{\mathrm M(a^{\mathrm n},b^{\mathrm n})}{2^n\mathrm M(a^{\mathrm n},c^{\mathrm n})}=\log y
\]
Thus, by the equations in Article 12, the limit of
$\frac12y^{-2^{n-2}}\cdot\surd{\frac{c^{\mathrm n}}{\mathrm M(a^{\mathrm n},b^{\mathrm n})}}$ for steadily increasing $n$, or, equivalently, the limit of $\frac12y^{-\frac14}r(y)$,
where $r(y)$ is written for $\surd{\frac{c}{\mathrm M(a,b)}}$, for positive $y^{\frac14}$ decreasing to zero, equals unity.

\Needspace{27\baselineskip}
If we further denote $\surd{\frac{a}{\mathrm M(a,b)}}$ and $\surd{\frac{b}{\mathrm M(a,b)}}$ by $p(y)$ and $q(y)$, then from
$aa=bb+cc$ and the two equations for $\surd{\mathrm M(a,b)}$ in Article 12 it follows that
\Needspace{10\baselineskip}
\[
\begin{aligned}p(y)&=1+r(y^4)+r(y^{16})+r(y^{64})\\&\quad+.\;.+r(y^{2^{2n}})+.\;.\end{aligned}
\]
\Needspace{10\baselineskip}
\[
\begin{aligned}q(y)&=1-r(y^4)+r(y^{16})+r(y^{64})\\&\quad+.\;.+r(y^{2^{2n}})+.\;.\end{aligned}
\]
\Needspace{10\baselineskip}
\[
\begin{aligned}r(y)^4&=p(y)^4-q(y)^4\end{aligned}
\]
\Needspace{25\baselineskip}
The series for $p,q,r$ that proceed in integral powers of $y^{\frac14}$ and satisfy these conditions are found, as far as the expansion is carried out, to have the form:
\Needspace{10\baselineskip}
\[
\begin{aligned}p(y)&=1+2y+2y^4+2y^9\\&\quad+.\;.+2y^{nn}+.\;.\;.\end{aligned}
\]
\Needspace{10\baselineskip}
\[
\begin{aligned}q(y)&=1-2y+2y^4-2y^9\\&\quad+.\;.\pm2y^{nn}\mp.\;.\;.\end{aligned}
\]
\Needspace{10\baselineskip}
\[
\begin{aligned}r(y)&=2y^{\frac14}+2y^{\frac94}+2y^{\frac{25}4}+2y^{\frac{49}4}\\&\quad+.\;.+2y^{(n+\frac12)^2}+.\;.\;.\end{aligned}
\]

\Needspace{29\baselineskip}
\SourceLeaf{393}%
That the law indicated here for forming the terms is generally valid seems to have been proved by \textsc{Gauss}, among other ways, as follows. Alongside the corresponding equations referring to series with two arguments, which will appear below in Articles 23 and 25, \textsc{Gauss} made the following note:]

\begingroup\itshape
On the theory of decomposing numbers into four squares.

The theorem: the product of two sums of four squares is itself a sum of four squares, is expressed most simply as follows:

let $l,m,\lambda,\mu,\lambda',\mu'$ be six complex numbers such that $\lambda,\lambda'$ and $\mu,\mu'$ are conjugate pairs.
Let $\mathrm N$ denote the norm. Then
\Needspace{10\baselineskip}
\[
(\mathrm Nl+\mathrm Nm)(\mathrm N\lambda+\mathrm N\mu)
=\mathrm N(l\lambda+m\mu)+\mathrm N(l\mu'-m\lambda')
\]
\par\endgroup
\Needspace{20\baselineskip}
[and therefore also
\Needspace{18\baselineskip}
\[
\begin{aligned}
&\{\mathrm N(n+in_{\prime})+\mathrm N(n_{\prime\prime}+in_{\prime\prime\prime})\}
\\&\quad\{\mathrm N(1-i)+\mathrm N(1+i)\}\\
&\qquad=\mathrm N\{(n+n_{\prime}+n_{\prime\prime}-n_{\prime\prime\prime})
\\&\qquad\quad+i(-n+n_{\prime}+n_{\prime\prime}+n_{\prime\prime\prime})\}\\
&\qquad\quad+\mathrm N\{(n+n_{\prime}-n_{\prime\prime}+n_{\prime\prime\prime})
\\&\qquad\quad-i(+n-n_{\prime}+n_{\prime\prime}+n_{\prime\prime\prime})\}
\end{aligned}
\]
From this the following two theorems can easily be derived. In them, different representations of a number as a sum of four square numbers refer to the different systems of values of the four roots, taking account of both the signs and the order of the roots; further, zero is included among the even numbers.

\Needspace{10\baselineskip}
If four times a number of the form $4k+1$ can be represented by four odd square numbers, then the number itself can be represented half as often by one odd and three even square numbers; conversely, if a number can be represented in the latter way, then four times that number can be represented twice as often in the former way.

\Needspace{10\baselineskip}
If four times a number of the form $4k+3$ can be represented by four odd square numbers, then the number itself can be represented half as often by one even and three odd square numbers; conversely, if a number can be represented in the latter way, then four times that number can be represented twice as often in the former way.

\Needspace{21\baselineskip}
With the aid of these two theorems it can be proved directly that the series above satisfy the equation $r(y)^4=p(y)^4-q(y)^4$, and further, by \SourceLeaf{394}%
repeated application of this relation, that if, according to the scheme
\Needspace{10\baselineskip}
\[
p+q=2p'',\quad p-q=2r'',\quad(p'')^4-(r'')^4=(q'')^4
\]
the terms of the algorithm for an arithmetic–geometric mean with positive even indices are formed from the values of the squares of the first two series, namely $pp$ and $qq$, as initial terms, then also
\Needspace{10\baselineskip}
\[
p^{2\mathrm n}=p(y^{2^{2n}}),\quad q^{2\mathrm n}=q(y^{2^{2n}}),\quad r^{2\mathrm n}=r(y^{2^{2n}})
\]
holds. Passing to the limit of $n$ therefore gives, by Article 12:
\Needspace{10\baselineskip}
\[
\mathrm M(pp,qq)=1,\qquad
\frac\pi2\frac{\mathrm M(pp,qq)}{\mathrm M(pp,rr)}=-\tfrac12\log y\ ]
\]

\clearpage
\endgroup
